\begin{solapartat}
\punts{0,75} Es tracta d'una desintegració $\beta^-$, per tant, s'emeten un electró i un antineutrí.

A més, la reacció ha de complir la conservació del nombre màssic nuclear i, a més, ha de conservar la càrrega; amb aquestes premisses, la reacció de desintegració serà:
\[ {}^{137}_{\phantom{1}55}\mathrm{Cs} \to {}^{x}_{y}\mathrm{Ba} + {}^{\phantom{-}0}_{-1}e + \bar{\nu}_e \]
Per una banda, es compleix la conservació del nombre màssic:
\[ 137 = x + 0 \;\Rightarrow\; x = 137 \]
I per a la conservació de càrrega:
\[ 55 = y - 1 \;\Rightarrow\; y = 56 \]
Per tant, la reacció completa és:
\[ {}^{137}_{\phantom{1}55}\mathrm{Cs} \to {}^{137}_{\phantom{1}56}\mathrm{Ba} + {}^{\phantom{-}0}_{-1}e + \bar{\nu}_e \]
\destaca{Alternativament,} es pot escriure ${}^{\phantom{-}0}_{-1}\beta$ o $e^-$ enlloc de ${}^{\phantom{-}0}_{-1}e$.

Si no s'indica l'antineutrí, cal restar \destaca{0,25 p}.

No cal que l'estudiant detalli les operacions relacionades amb la conservació de la càrrega i del nombre màssic, són operacions molt senzilles que es poden fer de cap; per tant, només es valorarà si ha estat capaç de deduir correctament els diferents termes de la reacció.

\medskip
\punts{0,25} El nombre de nuclis al cap d'un temps ve donada per:
\[ N = N_0\,e^{-\lambda t} \]
El coeficient de desintegració $\lambda$ s'obté del temps de vida mitja: $\lambda = \frac{1}{t_{vm}} = 1,05\times 10^{-9}\ \mathrm{s^{-1}}$

\punts{0,25} i per tant l'activitat al cap de 40 anys és:
\[ t = 40\ \text{anys} = 1,26\times 10^{9}\ \mathrm{s} \]
\[ N(40\ \text{anys}) = 100\%\,e^{-\lambda t} = 100\%\,e^{-1,05\times 10^{-9}\times 1,26\times 10^{9}} = 26,6\ \% \]
\destaca{Alternativament.} Atès que els estudiants estan més familiaritzats amb el temps de semidesintegració, s'han considerat correctes aquelles resolucions que estiguin ben plantejades tant si l'estudiant interpreta que el paràmetre proporcionat a l'enunciat és la vida mitjana com si interpreta que és el període de semidesintegració. Així, també es considera correcte el següent plantejament:

el coeficient de desintegració $\lambda$ s'obté del temps de semidesintegració:

$N(t) = \frac{N_0}{2} = N_0\,e^{-\lambda t_{1/2}}$ i per tant, $t_{1/2} = \frac{\ln 2}{\lambda}$.

I obtenim:
\[ t_{1/2} = 30,17\times 365\times 24\times 3600 = 9,51\times 10^{8}\ \mathrm{s} \]
\[ \lambda = \frac{\ln 2}{t_{1/2}} = \frac{\ln 2}{9,51\times 10^{8}} = 7,28\times 10^{-10}\ \mathrm{s^{-1}} \]
\punts{0,25} i per tant l'activitat al cap de 40 anys és:
\[ t = 40\ \text{anys} = 1,26\times 10^{9}\ \mathrm{s} \]
\[ N(40\ \text{anys}) = 100\%\,e^{-\lambda t} = 100\%\,e^{-7,28\times 10^{-10}\times 1,26\times 10^{9}} = 39,9\ \% \]
No cal que l'estudiant faci els càlculs dins el SI; per exemple, si ha operat amb anys també es considerarà correcte.
\end{solapartat}

\begin{solapartat}
\punts{0,5} L'activitat al cap d'un temps ve donada per:
\[ A = A_0\,e^{-\lambda t} \]
Si aïllem el temps:
\[ \frac{A}{A_0} = e^{-\lambda t} \Longrightarrow \ln\left(\frac{A}{A_0}\right) = -\lambda t \Longrightarrow t = -\frac{1}{\lambda}\ln\left(\frac{A}{A_0}\right) = \frac{1}{\lambda}\ln\left(\frac{A_0}{A}\right) \]
\punts{0,5} I si substituïm:
% DUBTE: la pauta només presenta el càlcul amb lambda = 7,28e-10 s^-1 (interpretant 30,17 anys com a
% període de semidesintegració); amb la vida mitjana (lambda = 1,05e-9 s^-1) sortirien uns 84 anys.
\[ t = \frac{1}{\lambda}\ln\left(\frac{A_0}{A}\right) = \frac{1}{7,28\times 10^{-10}}\ln\left(\frac{6,00\times 10^{5}}{3,70\times 10^{4}}\right) = 3,82\times 10^{9}\ \mathrm{s} = 121\ \text{anys} \]
\punts{0,25} Per tant, hauran de passar aproximadament 121 anys des de l'accident perquè la zona sigui habitable, és a dir, fins aproximadament el 2107.
\end{solapartat}
